% 第8章第7节：本章小结。

\section{本章小结}
\label{sec:ode-summary}

本章从“只知道每个瞬间的变化率，能否确定整条轨迹”出发。第~\ref{sec:ode-ivp}节把问题写成初值问题：向量场规定瞬时方向，积分方程把方向累计成状态；输入只在时间上可测时，解放宽为几乎处处满足方程的绝对连续曲线，而不是任意连续曲线。第~\ref{sec:ode-well-posedness}节用同一个积分方程回答了章首的第一组问题：在关于状态局部Lipschitz的条件下，积分映射在短时间的曲线空间里是压缩映射，解局部存在且唯一；局部解按唯一性拼接成最大解，它在有限时间内停止的唯一方式是离开状态集合中的每个紧集，线性增长条件借助Grönwall不等式排除了这种情形；初值、输入与模型的误差在有限时间内至多按$e^{Lt}$放大，可微时由变分方程给出一阶灵敏度。只有连续性时解可能不唯一，增长过快时解可能在有限时间爆炸，这两个反例标出了条件的边界。

第~\ref{sec:ode-analytic-solutions}节在可以求解的方程上落实这些结论。可分离变量方程把时间写成路程除以速度的积分，统一解释了爆炸与有限时间离开平衡点；常数变易公式把解分成初值的记忆与输入的累计；矩阵指数与状态转移矩阵给出线性系统的完整解法，其中时变系数的传播必须保留时间次序。零阶保持离散化在输入分段常值时是精确的，这一点把连续模型与离散模型准确地联系起来。

第~\ref{sec:ode-stability}节回答长期行为。线性系统的稳定性由特征值相对虚轴的位置决定，但非正规矩阵可能在衰减之前先放大；Lyapunov方程为稳定的线性系统提供一个单调下降的二次度量，线性化把这一结论推广到非线性平衡点的邻域。要对更大的区域下结论，需要Lyapunov函数；能量只是非严格下降时，LaSalle不变性原理用$\omega$极限集的不变性排除那些不能持续停留的零导数点。梯度流说明下降有下界并不等于状态收敛，极限环与分岔说明长期行为不只是趋于平衡点。持续输入下，稳定系统会遗忘初态，却会把输入的大小留在状态的最终范围中。

第~\ref{sec:ode-numerics}节与第~\ref{sec:ode-hamiltonian}节回答计算问题。一步方法的全局误差由相容性与稳定性共同控制：局部截断误差$O(h^{p+1})$经约$T/h$步累计为$O(h^p)$，舍入与隐式方程的残差则随步数增加而积累。刚性问题中，显式方法的步长受稳定域而非精度限制，A稳定的隐式方法以每步求解方程为代价解除这一限制，单侧Lipschitz条件说明了它们的误差界为何不随刚性恶化。自适应步长控制局部误差，却不能保证全局误差或长期的结构性质。对保守系统，精确流保持能量、辛结构与相空间体积，三者互不蕴含；辛Euler与Leapfrog保持辛结构而只近似保持能量，在合适的步长下能量误差有界而不漂移，这是阶数更高的通用方法做不到的。

两个贯穿例子概括了向量场结构的作用。温度模型和阻尼摆是耗散的：误差与初值记忆随时间消退，长期行为由稳定性理论刻画，数值方法要关心的是稳定域与刚性。无阻尼摆是保守的：误差不会自动消退，长期行为是在能量面上的永恒运动，数值方法要关心的是能否尊重这种结构。区分这些结论时，还要注意每一条结论作用于真实轨迹还是离散近似，以及它在有限时间内还是在长期成立。

本章的状态是有限维向量。若未知量是一根杆上的温度分布或一张膜的位移，状态就成为空间区域上的函数，方程中出现对空间变量的导数，边界上的条件也成为问题的一部分。第~\ref{chap:partial-differential-equations}章研究这种情形，本章的时间推进方法将在那里作为空间离散之后的求解工具再次出现。
